\section{A common sampling schedule}\label{sec:schedule}

The transfer theorem explains what a useful schedule must achieve.
At each checkpoint, the joint residue entropy must be almost maximal,
while the entropy of repeated continuations must be small compared
with that joint entropy. A much more accurate marginal estimate is
needed only at the last checkpoint, where it starts the backward
coverage induction. We now construct one schedule that supplies these
estimates for every prime batch.

The same construction also supplies the inputs needed after coverage
has been established: a persistent joint entropy of order \(T\) for
each batch near \(T\), a smaller total alphabet budget for all preceding
batches, and one common terminal-time law whose small translates are
close in total variation. These requirements have different roles.
The displacement bound pays for information revealed in the coverage
argument; the final translation bound will permit the information costs
of different batches to be added.

Throughout the section, the walk is an arbitrary fixed infinite
self-avoiding \(D\)-step walk. No zero-avoidance assumption is used.
We first give the actual schedule and prove its entropy and displacement
estimates, then apply Theorem~\ref{thm:transfer} to obtain coverage.

\subsection{Choice of constants and scales}

Fix \(D\ge1\), and set
\[
\beta_0=\frac1{200},\qquad e_0=\frac{\beta_0}{50}.
\]
We will first choose a fixed iteration count \(l_0\) and a sufficiently
small constant \(0<g_0<10^{-2}\), depending only on \(D,e_0\).
They determine a constant \(c_0>0\); only then do we choose a sufficiently
large fixed \(B>2\). The precise requirements are specified below.
All these constants are fixed before the asymptotic parameter \(M\),
and are independent of the walk and starting time.

Let \(M\) be a positive integer tending to infinity, and put
\[
X_m=100^m\quad(\lceil M/2\rceil\le m\le M).
\]
Use every batch \(\mathcal B_T\) with
\begin{equation}\label{eq:batches}
T=B^j,\quad j\in\Z,\qquad
X_m\le\log T\le1.05X_m
\quad\text{for some }\lceil M/2\rceil\le m\le M.
\end{equation}
The batches are disjoint because \(B>2\).
Write \(T_+=\exp(1.05X_M)\).
For this value of \(M\), the union of all these batches is the finite
prime selection that will be used in Theorem~\ref{thm:sieve}.

Starting at time \(0\) on a fixed infinite self-avoiding \(D\)-step
walk, apply the transition of Lemma~\ref{lem:enrichment} with scales
\(U\) in decreasing order. For each \(m\), use the three bands
in Table~\ref{tab:bands}. In a band with parameter \(g\), start at its
upper endpoint in \(\log U\), and successively subtract
\(2\log(1/g)\) as long as the resulting scale stays in that band.
Equivalently, successive \(U\)'s have ratio \(g^2\).
All randomizations are fresh conditional on the current exact walk time.
No prime subset or sign choice enters these transitions.

\begin{table}[ht]
\centering
\begin{tabular}{lll}
\toprule
Band & Interval for \(\log U\) & Transition parameter\\
\midrule
Top & \([0.8X_m,1.3X_m]\) & \(g_0\)\\
Middle & \([0.20X_m,0.30X_m]\) & \(g_1=e^{-M^3}\)\\
Bottom & \([0.05X_m,0.10X_m]\) & \(g_2=e^{-100M^{20}}\)\\
\bottomrule
\end{tabular}
\caption{The three bands in each scale window. Windows are processed
in decreasing \(m\), and scales decrease within each band.}
\label{tab:bands}
\end{table}

After all transitions in all windows have been performed, add one
independent uniform offset in \(\{0,\ldots,N-1\}\), where
\(N=\lceil T_+^{10}\rceil\).
Denote the resulting random time by \(t_*\).
Thus all batches will be measured at the same time law.

A \emph{checkpoint} is a stage just before a specified transition,
or the terminal stage. It also denotes the random walk time at that
stage. Conditioning on an exact start at a checkpoint means fixing that
walk time and then running the remaining schedule with fresh randomness.
Such conditional experiments are defined at every nonnegative integer
start, whether or not it has positive probability in the experiment
starting from time \(0\).
The term ``exact start'' specifies a new continuation experiment,
not conditioning on an event that may have probability zero.
\begin{lemma}[Future shifts and smoothing]\label{lem:future-shift}
Let \(V\) be a scheduled scale in a band with parameter \(g\).
For all sufficiently large \(M\), the maximum total forward time shift
from the checkpoint before \(V\) to any later checkpoint, including
\(t_*\), is bounded by a deterministic number \(S_V\) satisfying
\[
\log(1+S_V)\le2gV.
\]
Consequently the entropy of its physical displacement is \(O_D(gV)\),
for every starting law.
Furthermore, conditional on any earlier exact checkpoint start,
\begin{equation}\label{eq:smoothing}
\TV\bigl(\mathcal L(t_*),\mathcal L(t_*+a)\bigr)
\le T_+^{-9}\qquad
(0\le a\le T_+,\ a\in\Z).
\end{equation}
\end{lemma}
\begin{proof}
One may take the deterministic bound
\[
 S_V=N-1+\sum_{U\text{ remaining}}\lfloor\exp(g_U U)\rfloor,
\]
where the sum includes the transition at $V$ and $g_U$ is the parameter
of the band containing $U$. Every summand and the offset length $N$
are fixed by the schedule, before any exact start is chosen. This bounds
the total advance from every such start.
The number of operations is \(O_{g_0}(X_M)=\exp(O_{g_0}(M))\).
A transition at scale \(U\) shifts time by at most \(\exp(gU)\).
Within the same band all future log-lengths are at most \(gV\).
Between successive bands the decrease in \(U\) is at least
\(\exp(cX_m)\) for an absolute \(c>0\), including the passage from
the bottom band for \(m\) to the top band for \(m-1\):
\(0.05X_m-1.3X_{m-1}=0.037X_m\).
This decrease dominates every ratio of transition parameters, because
\(\log(1/g_2)=100M^{20}=o(X_{\lceil M/2\rceil})\).
Thus later bands also have log-length at most \(gV\) for large \(M\).

The logarithm of the final offset length is \(O(X_M)\).
On the other hand, for the least possible scale and parameter,
\[
\log(gV)\ge 0.05X_{\lceil M/2\rceil}-100M^{20},
\]
so \(X_M=o(gV)\), uniformly.
Taking the logarithm of the sum of all remaining lengths adds only
\(O_{g_0}(M)\). The bound \(2gV\) follows, and
Lemma~\ref{lem:displacement} gives the entropy assertion.

Write \(t_*=\widetilde t+U_N\), where \(U_N\) is the final uniform
offset, independent of all earlier choices.
A translate by \(a\le T_+\) changes its law by at most \(a/N\).
Mixtures and deterministic translations do not increase this bound,
including after an earlier exact start is fixed.
This proves \eqref{eq:smoothing}.
\end{proof}

\subsection{Iterating and preserving joint entropy}

For a fixed batch write \(f_Z(s)\) as in Section~\ref{sec:geometry}.
Every use of Lemma~\ref{lem:enrichment} below has \(q\le k_T\).
For each entropy estimate we begin at a scale where
\(q\le k_T\); all subsequent scales in that block are smaller.
The other limiting conditions hold uniformly:
\[
\log U\ge cX_{\lceil M/2\rceil},\qquad
\log(1/g)=O(M^{20}),\qquad
\log\log T=O(M).
\]
In particular \(g^CU/\log T\to\infty\) for every fixed \(C>0\).

Suppose \(l\) successive transitions in one band end just before its
next scale \(V\), and that the initial subset size is at most \(k_T\).
The output size of one transition and the input size of the next are
both exactly \(\lfloor g^2U/L_*\rfloor\).
Writing the normalized entropy deficit as \(d_j\),
Equation~\eqref{eq:enrichment} gives
\(d_{j+1}\le d_j/2+C_Dg\). Since \(d_0\le1\), iteration gives
\(d_l\le2^{-l}+2C_Dg\), and hence
\begin{equation}\label{eq:iterate}
\frac{f_Z(\lfloor V/L_*\rfloor)}
     {\lfloor V/L_*\rfloor L_*}
\ge1-2^{-l}-O_D(g)
\end{equation}
at that checkpoint.
The same bound, with another \(O_D(g)\) error, holds at every later
checkpoint. Indeed Lemmas~\ref{lem:future-shift}
and~\ref{lem:displacement} show that the entropy loss at this fixed
subset size is \(O_D(gV)\), whereas
\(\lfloor V/L_*\rfloor L_*\ge V/2\) for large \(M\).
The argument is uniform for every initial law, including an exact
start fixed before these \(l\) transitions. The later entropies here
are marginal entropies in that experiment; we are not conditioning on
the exact later time, which would determine its walk position.

Choose a small absolute \(a_0>0\) such that
\(a_0T/L_*\le k_T\) for large \(T\), using \eqref{eq:prime-count}.
For a batch \(T\) in window \(m\), begin an entropy estimate at the
first top-band scale \(U\le a_0T\).
Its size lies between \(g_0^2a_0T\) and \(a_0T\).
Choose a fixed integer \(l_0\) large enough that
\(2^{-l_0}<e_0/3\), and take \(g_0\) small enough to make the total
error in \eqref{eq:iterate} and its preservation less than \(2e_0/3\).
These choices can be made before \(B\) and \(M\).
Run \(l_0\) transitions and let \(q_T=\lfloor V/L_*\rfloor\) be the
recorded size. All these scales lie in the top band for large \(M\).
At every later checkpoint, for its point variable \(Z\),
\begin{equation}\label{eq:global-entropy}
f_Z(q_T)\ge(1-e_0)q_TL_*,
\qquad q_TL_*\ge c_0T
\end{equation}
with \(c_0>0\) depending only on the fixed choices.
For example any sufficiently small constant below
\(a_0g_0^{2(l_0+1)}/2\) works in the second inequality.

By \eqref{eq:prime-count}, \(k_{T'}\log(2T')\le C T'\).
Since the batches are a subsequence of the geometric progression
\(B^j\), we can now choose \(B>2\) so large that for every used batch
\begin{equation}\label{eq:smaller-alphabet}
\sum_{T'<T}k_{T'}\log(2T')\le e_0c_0T.
\end{equation}
In fact the left side is at most \(CT/(B-1)\).
This completes the fixed-constant choices in the order
\(D,e_0\), then \(l_0,g_0,c_0\), then \(B\), and finally \(M\).
The top band has now supplied a batch entropy of order \(T\) that
survives the remaining schedule, and the geometric spacing makes all
smaller residue alphabets inexpensive compared with it. We next build
the more accurate conditional estimates needed for coverage.

\subsection{Intermediate checkpoints and almost uniform marginals}

These checkpoints will transmit coverage backward from an accurate
one-coordinate base case. Their spacing must allow joint entropy to
grow while keeping subsequent displacement data inexpensive.
For each batch \(T\), put
\begin{equation}\label{eq:checkpoint-parameters}
\beta_i=2^{-i}\beta_0,\qquad r_0=\tfrac14,\qquad
r_{i+1}=r_i-2\beta_i.
\end{equation}
Stop at the least \(J=J(T)\) with \(\beta_J\log T\le M^{20}\).
Since \(\beta_0\log T>M^{20}\) for sufficiently large \(M\), we have
\(J\ge1\); minimality then gives
\begin{equation}\label{eq:checkpoint-ranges}
J=O(M),\qquad
M^{20}/2<\beta_J\log T\le M^{20},\qquad
0.23<r_i\le0.25.
\end{equation}
The last assertion follows from
\(r_i=1/4-4\beta_0(1-2^{-i})\).
Let \(A_i\) be the checkpoint just before the first scheduled operation
with \(U\le T^{r_i}\).
All these checkpoints lie in the middle band for this \(m\), since
\[
0.23X_m<r_i\log T\le0.2625X_m.
\]
Its scale is in \((g_1^2T^{r_i},T^{r_i}]\).
Let \(S_i\) be a deterministic upper bound for the remaining time
shift from \(A_i\) supplied by Lemma~\ref{lem:future-shift}, and define
\[
 \mathcal D_i=\{v\in\Z[i]:|v|\le DS_i\}.
\]
This set does not depend on the exact start. From every exact start
\(A_i=a\), every possible displacement to any later checkpoint,
including \(z_{t_*}-z_a\), belongs to \(\mathcal D_i\). Since the scale
at \(A_i\) is at most \(T^{r_i}\), lattice-point counting gives
\begin{equation}\label{eq:native-data}
 \log|\mathcal D_i|\le C_D T^{r_i}.
\end{equation}
Consequently, uniformly in the starting law, for any later checkpoint
\(A\) we have
\begin{equation}\label{eq:checkpoint-displacement}
 H(z_A-z_{A_i})\le C_D T^{r_i}.
\end{equation}
The common support bound is stronger than a separate conditional entropy
bound at each exact start. It will control the unconditional entropy
of displacement data when the next checkpoint is random.

Fix an exact start at \(A_i\), where \(i<J\), and let
\(Y=z_{A_{i+1}}\).
There is a deterministic subset size \(s_i\), depending on the batch
and the schedule but not on that start, such that
\begin{equation}\label{eq:checkpoint-entropy}
f_Y(s_i)\ge(1-e^{-cM^3})s_iL_*,
\qquad
s_iL_*\ge\exp(r_i\log T-O(M^6))
\end{equation}
for a fixed \(c>0\).
To see this, take \(\lceil M^3\rceil\) iterations starting at \(A_i\).
The first log-scale differs from \(r_i\log T\) by less than \(2M^3\),
and these iterations consume \(O(M^6)\) log-scale.
The starting subset size is at most \(k_T\), since
\(U\le T^{r_i}\le T^{1/4}=o(T)\) and
\(k_TL_*\asymp T\). They fit before \(A_{i+1}\), because
\[
(r_i-r_{i+1})\log T=2\beta_i\log T>2M^{20}.
\]
Equation~\eqref{eq:iterate}, preservation, and \(g_1=e^{-M^3}\)
prove the first bound, and the recorded scale proves the second.

Finally, conditional on any exact start at \(A_J\), proceed to the
bottom band for this \(m\).
Take \(\lceil200M^{20}\rceil\) iterations, recording a subset size \(s\).
Begin at its first scale, which is at most
\(\exp(0.10X_m)=o(T)\), so again the initial subset size is at most
\(k_T\). The required log-scale span is \(O(M^{40})=o(X_m)\), so the
iterations and their recording scale fit in that band. Their deficit, including preservation to \(t_*\),
is at most
\[
2^{-\lceil200M^{20}\rceil}+O_D(e^{-100M^{20}})
\le e^{-90M^{20}}.
\]
Concavity of \(f\) gives \(f(1)\ge f(s)/s\).
Thus in every such exact-start experiment,
\begin{equation}\label{eq:terminal-entropy}
f_{z_{t_*}}(1)\ge(1-e^{-90M^{20}})L_*.
\end{equation}
Equations~\eqref{eq:checkpoint-entropy} and~\eqref{eq:terminal-entropy}
are the two inputs to the coverage induction: high joint entropy
between checkpoints and an extremely accurate one-coordinate base case.
Figure~\ref{fig:schedule} summarizes the two orders now in play:
sampling moves forward through the schedule, while coverage will be
proved by induction from $A_J$ back to $A_0$.

\begin{figure}[ht]
\centering
\begin{tikzpicture}[
  >=Latex,font=\small,
  stage/.style={draw,align=center,inner sep=4pt,minimum height=1.05cm},
  proof/.style={draw,align=center,inner sep=4pt,minimum height=0.85cm}]
\node[anchor=west] at (-1.2,1.05)
  {Forward sampling from one window onward};
\node[stage,text width=2.1cm] (top) at (0,0)
  {Top band\\joint entropy};
\node[stage,text width=2.6cm] (mid) at (3.0,0)
  {Middle band\\\(A_0,\ldots,A_J\)};
\node[stage,text width=2.4cm] (bottom) at (6.0,0)
  {Bottom band\\marginals};
\node[stage,text width=3.4cm] (rest) at (9.4,0)
  {Smaller windows\\then uniform offset};
\node[stage,text width=0.8cm] (terminal) at (12.6,0)
  {\(t_*\)};
\draw[->] (top.east) -- (mid.west);
\draw[->] (mid.east) -- (bottom.west);
\draw[->] (bottom.east) -- (rest.west);
\draw[->] (rest.east) -- (terminal.west);

\node[anchor=west] at (-1.2,-1.05)
  {Backward coverage induction for exact-start continuation laws};
\node[proof,text width=2.8cm] (base) at (1.15,-2.02)
  {Coverage at \(A_J\)\\base case};
\node[proof,text width=2.8cm] (step) at (5.65,-2.02)
  {\(A_{i+1}\) implies \(A_i\)\\\(i=J-1,\ldots,0\)};
\node[proof,text width=2.8cm] (goal) at (10.65,-2.02)
  {Coverage at \(A_0\)};
\draw[->] (base.east) -- (step.west);
\draw[->] (step.east) -- (goal.west);
\end{tikzpicture}
\caption{Sampling order and proof order have separate rows. The top row
suppresses individual transitions and shows the common terminal time.
The bottom row proves coverage of the terminal law from every exact
start at each checkpoint. In the step from \(A_{i+1}\) to \(A_i\), fix
the exact start at \(A_i\), run forward to \(A_{i+1}\), then sample
independent continuations conditional on that exact next time.
The induction reverses the checkpoint index; all sampled times move
forward along the fixed walk.}
\label{fig:schedule}
\end{figure}

\subsection{Coverage from every exact checkpoint}\label{sec:checkpoints}
\label{sec:coverage}

We now apply the earlier transfer theorem to the estimates just proved.
The coordinate groups are \(\Z[i]/(\pi)\), their maps are reduction
modulo the two factors over each prime of \(\mathcal B_T\), and the
abstract average \(\overline L\) is exactly \(L_*\).
Averages over signed primes assign equal weight to all \(2k_T\)
prime-factor pairs.

Recall that \((\tau,\beta)\) coverage means that a law on an alphabet
of size \(p\) assigns mass at least \(\tau/p\) outside fewer than
\(p^{1-\beta}\) residues. We use the following floors and exceptional
fractions at the checkpoints already constructed:
\begin{align}
 \tau_0&=\tfrac14,& \delta_i&=2^{-i-6},&
 \tau_{i+1}&=\tau_i+4\delta_i,
 \label{eq:coverage-parameters}\\
 \eta_0&=\tfrac1{100},&&&
 \eta_{i+1}&=\eta_i\delta_i/100.\label{eq:native-eta}
\end{align}
Here \(\beta_i\), \(r_i\), and \(J\) are those of
\eqref{eq:checkpoint-parameters}--\eqref{eq:checkpoint-ranges}.
In particular,
\[
 \tau_i=\frac38-2^{-i-3}<\frac38,\qquad
 \eta_i=100^{-(i+1)}2^{-i(i-1)/2-6i}.
\]
Since \(J=O(M)\), these formulas imply, uniformly over \(0\le i\le J\),
\begin{equation}\label{eq:coverage-parameter-bounds}
 \log(1/\delta_i)=O(M),\qquad
 \eta_i\ge\exp(-O(M^2)),
 \qquad \eta_i\delta_i\beta_{i+1}\ge\exp(-O(M^2)).
\end{equation}
The last estimate is needed only for \(i<J\).

For a specified exact start \(a\) at checkpoint \(A_i\), denote the
terminal residue law of its continuation experiment by
\[
 b_{i,a,\pi}(u)=
 \Pr\{z_{t_*}\bmod\pi=u:
       \text{the remaining schedule is started at }A_i=a\}.
\]
The colon specifies the experiment, including at starts that have
probability zero under the schedule begun at time \(0\).

\begin{proposition}[Terminal residue coverage]\label{prop:coverage}
For all sufficiently large \(M\), uniformly over the batches,
\(0\le i\le J\), and all exact start times \(a\ge0\) at \(A_i\),
all but a fraction at most \(\eta_i\) of the signed primes satisfy
\begin{equation}\label{eq:coverage}
 \#\{u\in\Z[i]/(\pi):b_{i,a,\pi}(u)<\tau_i/p\}
 <p^{1-\beta_i}.
\end{equation}
\end{proposition}

\begin{proof}
We begin at \(A_J\) using the particularly accurate marginal estimate,
then verify every numerical hypothesis of Theorem~\ref{thm:transfer}
for the passage from \(A_{i+1}\) to \(A_i\).

\paragraph{The terminal base.}
Fix any exact start at \(A_J\).
By \eqref{eq:terminal-entropy}, the average relative entropy of the
terminal residue laws from their respective uniform laws is at most
\(e^{-90M^{20}}L_*\).
If a signed prime fails \eqref{eq:coverage}, its low-probability set
has at least \(p^{1-\beta_J}\) elements. The uniform law assigns that
set at least \((1-\tau_J)p^{-\beta_J}\) more mass than its terminal law.
Thus the total variation distance is at least
\[
 (1-\tau_J)p^{-\beta_J}
 \ge c e^{-M^{20}},
\]
using \(p\le2T\), \(\tau_J<3/8\), and
\(\beta_J\log T\le M^{20}\).
Pinsker's inequality in natural logarithms
\cite[Lemma 11.6.1]{CoverThomas} makes each such entropy deficit at least
\(c'e^{-2M^{20}}\).
Averaging the lower bound over the failing signed primes bounds their
fraction by
\[
 C L_*e^{-88M^{20}}
 =\exp(-88M^{20}+O(M))<\eta_J
\]
for all sufficiently large \(M\). This proves the base at every exact
start. The terminal marginal deficit has the accuracy needed to absorb
the power-saving exceptional-set threshold.

\paragraph{The shared data and its cost.}
Suppose coverage has been proved at \(i+1\), and fix an arbitrary
exact start \(A_i=a\). Run the schedule to the next checkpoint and put
\(A=A_{i+1}\) and \(Y=z_A\). Conditional on this exact time \(A\),
repeat the remaining continuation independently
\[
 n_i=\left\lceil T^{6\beta_i/5}\right\rceil
\]
times. Write \(h_1,\ldots,h_{n_i}\) for their physical displacements
from \(Y\), and \(\mathcal H=(h_1,\ldots,h_{n_i})\).
Every \(Y+h_\ell\) has the terminal law from the fixed start \(A_i=a\).
This one data vector is constructed before any auxiliary signed
coordinate subset is selected.

By \eqref{eq:native-data}, every \(h_\ell\) belongs to the same finite
set \(\mathcal D_{i+1}\), even as \(A\) varies. Hence its
\emph{unconditional} entropy in this fixed-start experiment satisfies
\[
 H(\mathcal H)\le n_i\log|\mathcal D_{i+1}|
 \le C_Dn_i T^{r_{i+1}}.
\]
This is not an assertion of independence before conditioning on \(A\).
Combining the support bound with \eqref{eq:checkpoint-entropy} yields
\begin{equation}\label{eq:coverage-data-cost}
 \kappa_i:=\frac{H(\mathcal H)}{s_iL_*}
 \le C_D\exp\left(-\frac45\beta_i\log T+O_D(M^6)\right)
 \le e^{-M^3}.
\end{equation}
Indeed \(n_i\le2T^{6\beta_i/5}\),
\(r_i-r_{i+1}=2\beta_i\), and
\(\beta_i\log T>M^{20}\) for \(i<J\).
The exponent \(4\beta_i/5\) is the margin left after paying for all
\(n_i\) continuations. By Lemma~\ref{lem:conditioning}, the joint
entropy bound and this common cost imply
\begin{equation}\label{eq:coverage-conditional-deficit}
 \E_{p,\pi}\bigl[\log p-H(Y\bmod\pi\mid\mathcal H)\bigr]
 \le \bigl(e^{-cM^3}+\kappa_i\bigr)L_*.
\end{equation}

\paragraph{Checking the transfer inequalities.}
In Theorem~\ref{thm:transfer} take
\[
 (\tau,\delta,\tau',\beta,\beta',\eta,\eta')
 =(\tau_i,\delta_i,\tau_{i+1},\beta_i,\beta_{i+1},
   \eta_i,\eta_{i+1}).
\]
The induction hypothesis supplies the conditional coverage at every
exact value of \(A\). Equation~\eqref{eq:checkpoint-entropy} supplies
its joint entropy input with \(\epsilon_i=e^{-cM^3}\).
The parameter estimates give
\[
 \epsilon_i+\kappa_i
 \le \eta_i\delta_i\beta_{i+1}/16,
 \qquad \eta_{i+1}=\eta_i\delta_i/100,
\]
for sufficiently large \(M\), which are the two entropy budget
requirements. For its concentration requirement, \(p\le2T\) gives
\[
 n_i p^{-\beta_i}
 \ge2^{-\beta_i}T^{\beta_i/5}
 \ge\frac12\exp(M^{20}/5).
\]
It follows that
\(\delta_i^2n_i p^{-\beta_i}/2\ge
\exp(M^{20}/5-O(M))\).
This quantity dominates both
\(\log p\le1.05\,100^M+\log2\) and
\(\log(100/\delta_i)=O(M)\). Consequently
\[
 p\exp(-\delta_i^2n_i p^{-\beta_i}/2)\le\delta_i/100.
\]
The two remaining inequalities are
\[
 \delta_i^{-1}<p^{\beta_{i+1}/2},\qquad
 \log(2e/\delta_i)\le(\beta_{i+1}/4)\log T.
\]
They follow from \(\log(1/\delta_i)=O(M)\) and
\(\beta_{i+1}\log T>M^{20}/2\).
All hypotheses of Theorem~\ref{thm:transfer} now hold. It proves
coverage at \(A_i=a\), and this start was arbitrary. Downward induction
establishes the proposition uniformly over the whole schedule.
\end{proof}

At \(A_0\), the conclusion gives mass at least \(1/(4p)\) outside
fewer than \(p^{199/200}\) residues for at least \(99\%\) of the signed
primes, from every exact start. The next section will use this coverage
and the common terminal-time law to test candidate starting residues
under zero avoidance. There the repeated tests will be independent
conditional on residue vectors, whose laws mix exact starts; that
conditioning will be treated separately.
